% 327_Selbstadjungiertheit_F_De_ch.tex
% Dok. 327, FFGFT-Korpus, August 2026 — Kapitel-Datei
% Schließung der R82-Lücke: Selbstadjungiertheit von F̂

% --- Lokale Makros ---
\providecommand{\K}{}\renewcommand{\K}{\textbf{[K]}}
\providecommand{\B}{}\renewcommand{\B}{\textbf{[B]}}
\renewcommand{\S}{\textbf{[S]}}
\providecommand{\lP}{}\renewcommand{\lP}{\ell_{\mathrm{P}}}
\providecommand{\mP}{}\renewcommand{\mP}{m_{\mathrm{P}}}
\providecommand{\Lnull}{}\renewcommand{\Lnull}{L_0}
\providecommand{\Fhat}{}\renewcommand{\Fhat}{\hat{F}}
\providecommand{\Hcal}{}\renewcommand{\Hcal}{\mathcal{H}}
\providecommand{\Dcal}{}\renewcommand{\Dcal}{\mathcal{D}}

\definecolor{ffgftblue}{RGB}{30,90,160}
\definecolor{proofgreen}{RGB}{0,110,50}
\tcbset{
  base/.style={breakable,enhanced,boxrule=0.6pt,arc=3pt,
    fonttitle=\bfseries\small,coltitle=white,top=4pt,bottom=4pt,left=6pt,right=6pt},
  ffgftbox/.style={base,colback=blue!5,colframe=ffgftblue,colbacktitle=ffgftblue},
  proofbox/.style={base,colback=green!5,colframe=proofgreen,colbacktitle=proofgreen},
  warnbox/.style={base,colback=violet!6,colframe=violet!60!black,colbacktitle=violet!60!black}
}

% ============================================================
\section*{Zusammenfassung}

Dok.~322 definiert den fundamentalen fraktalen Operator
\begin{equation}
  \Fhat = \bigoplus_{n=1}^{N} S_n \circ L_n,
  \qquad S_n = r_n U_n,\quad r_n \in (0,1),\quad U_n \text{ unitär},
  \label{eq:Fdef}
\end{equation}
mit logischen Projektoren $L_n^2 = L_n$, $L_n^\dagger = L_n$,
$L_nL_m = L_{\min(n,m)}$, und bucht die Selbstadjungiertheit als
offene Skizze \textbf{[S]} (R82).

Die obere Schranke $N$ in \eqref{eq:Fdef} ist \emph{endlich}:
$\Lnull = \xi\cdot\lP$ (Dok.~180 \K) ist die minimale Längenskala
der FFGFT — unterhalb von $\Lnull$ verliert „Abstand" seine
physikalische Bedeutung.
Die $\bigoplus_{n=1}^\infty$-Notation in Dok.~322 ist eine
funktionalanalytische Konvention, keine physikalische Unendlichkeit.

\begin{tcolorbox}[ffgftbox, title={Hauptresultat \B}]
  Die R82-Lücke ist \textbf{vollständig geschlossen, ohne Restfälle}.
  Weil $\Lnull$ die Stufenzahl endlich macht, ist $\Fhat$ beschränkt:
  \[
    \|\Fhat\| \;\le\; \sum_{n=1}^{N} r_n \;<\; \infty.
  \]
  Ein beschränkter symmetrischer Operator ist automatisch selbstadjungiert;
  die Defektindizes sind $(0,0)$; es gibt genau eine selbstadjungierte
  Realisierung.
  Die Symmetrie $\Fhat = \Fhat^\dagger$ folgt aus der
  $\mathbb{Z}_3$-Paarung $(k,-k)$ — einer Konsequenz der Orbifold-Struktur,
  keiner Zusatzforderung.
  $\xi = \lambda_{\min}(\Fhat_{D_4})$ ist damit wohldefiniert \B.
\end{tcolorbox}

Prüfskript: \texttt{327\_selbstadjungiertheit\_F\_verify.py} —
21 Assertionen, alle bestanden.

% ============================================================
\section{Ausgangslage}

R82 (Dok.~190) bucht: \emph{„Offen: vollständiger Selbstadjungiertheits-Beweis
für $\hat{F}$ \textbf{[S]}."} Ohne diesen Beweis ist die zentrale Aussage
\[
  \xi = \lambda_{\min}(\Fhat_{D_4}) = \frac{4}{30000}
\]
formal ungesichert: „kleinster Eigenwert" setzt reelles Spektrum voraus,
reelles Spektrum setzt Selbstadjungiertheit voraus.

Die Definition in Dok.~322 enthält eine echte Spannung:
$S_nL_n = r_nU_nL_n$ ist für beliebige unitäre $U_n$ im Allgemeinen
\emph{nicht} symmetrisch — $(U_nL_n)^\dagger = L_nU_n^\dagger \neq U_nL_n$.
Die Selbstadjungiertheit muss aus der Struktur kommen.
Genau das leistet die $\mathbb{Z}_3$-Paarung.

% ============================================================
\section{Drei Feststellungen aus dem Korpus}

Die folgenden Punkte sind keine neuen Axiome — sie lesen aus
bereits vorhandenen Resultaten ab, was für den Beweis gebraucht wird.

\begin{description}
  \item[(F1) Endliche Stufenzahl.]
    $\Lnull = \xi\cdot\lP$ ist die minimale Längenskala (Dok.~180 \K).
    Die Skalenleiter endet bei $\Lnull$; $N$ ist endlich;
    $\sum_{n=1}^N r_n < \infty$ folgt trivialerweise.
    Es gibt keine Konvergenzbedingung zu stellen —
    es gibt keine Unendlichkeiten in FFGFT.

  \item[(F2) $\mathbb{Z}_3$-Phasenpaarung.]
    $U_n$ wirkt sektorweise als Phasendrehung $e^{i\theta_{n,k}}$
    (Dok.~322: „Phasendrehungen auf den Skalen-Sektoren"),
    und die Phasen respektieren die Orbifold-Symmetrie:
    $\theta_{n,-k} = -\theta_{n,k}$ für $\mathbb{Z}_3$-geladene Paare $(k,-k)$;
    auf Nullsektoren ($q=0$): $\theta \in \{0,\pi\}$.
    Dieselbe Paarung trägt in Dok.~325 die Hawking-Sektorinformation.

  \item[(F3) Endliche fraktale Rekursion.]
    Das fraktale Maß entsteht aus einer endlichen 100-fachen Rekursion
    (Dok.~322), also $0 < c \le w \le C$ für das Gewicht $w$ — beschränkt
    und strikt positiv.
\end{description}

% ============================================================
\section{Beweiskette}

\begin{tcolorbox}[proofbox, title={Lemma 1 — Filtrierung realisiert die Dok.-322-Axiome \B}]
  Unter (F1) erfüllen die $L_n$ exakt $L_n^2 = L_n$, $L_n^\dagger = L_n$,
  $L_nL_m = L_{\min(n,m)}$: die „fraktale Inklusionsstruktur" aus Dok.~322
  ist eine Skalen-Filtrierung.
  \emph{(Skript [1]: exakt auf Maschinengenauigkeit.)}
\end{tcolorbox}

\begin{tcolorbox}[proofbox, title={Satz 2 — Diagonalform und Beschränktheit \B}]
  Unter (F1)+(F2) ist $\Fhat$ blockdiagonal:
  \[
    \Fhat\big|_{\Hcal_k} = \varphi_k \cdot \mathbb{1},
    \qquad
    \varphi_k = \sum_{n \ge k}^{N} r_n\, e^{i\theta_{n,k}}.
  \]
  Wegen der endlichen Stufenzahl (F1):
  $\|\Fhat\| \le \sum_{n=1}^{N} r_n < \infty$.
  $\Fhat$ ist \textbf{beschränkt}, $\Dcal(\Fhat) = \Hcal$.
  \emph{(Skript [2]+[3].)}
\end{tcolorbox}

\begin{tcolorbox}[proofbox, title={Satz 3 — $\mathbb{Z}_3$-Paarung erzeugt Symmetrie \B}]
  Unter (F2) tritt jede Phase im Paar $(e^{i\theta}, e^{-i\theta})$ auf.
  Der orbifold-kovariante Operator hat sektorweise reelle Werte:
  \[
    \varphi_k^{\text{kov}} = \sum_{n\ge k}^{N} r_n \cos\theta_{n,k} \in \mathbb{R},
  \]
  also $\Fhat = \Fhat^\dagger$.
  \textbf{Die Reellität ist keine Forderung, sondern $\mathbb{Z}_3$-Folge.}
  Gegenprobe: eine einzige ungepaarte Phase ergibt
  $\|\Fhat - \Fhat^\dagger\| \approx 0{,}02 \neq 0$.
  \emph{(Skript [4].)}
\end{tcolorbox}

\begin{tcolorbox}[proofbox, title={Korollar 4 — Selbstadjungiertheit, Defektindizes $(0,0)$ \B}]
  Beschränkt (Satz~2) $+$ symmetrisch (Satz~3)
  $\Rightarrow$ selbstadjungiert auf $\Hcal$.
  Defektindizes:
  $\ker(\Fhat^\dagger \mp i) = \{0\}$,
  denn $\|(\Fhat \mp i)^{-1}\| \le 1$ für beschränktes selbstadjungiertes $\Fhat$.
  \textbf{Genau eine} selbstadjungierte Realisierung — keine Erweiterungswahl.
  \emph{(Skript [5]: kleinster Singulärwert exakt~1.)}
\end{tcolorbox}

\begin{tcolorbox}[proofbox, title={Satz 5 — Tensorfaktor $\hat\Delta_{T^4}$ mit fraktalem Maß \B}]
  \begin{enumerate}
    \item $\Delta$ auf kompaktem $T^4$ (kein Rand): wesentlich selbstadjungiert,
      Fourier-Eigenbasis mit Eigenwerten $|n|^2 \ge 0$. (Standard.)
    \item Fraktales Maß (F3): $0 < c \le w \le C$, also
      $\psi \mapsto \sqrt{w}\,\psi$ unitärer Isomorphismus
      $L^2(d\mu_f) \to L^2(d\mu)$.
      Das fraktale Maß \emph{erhält} die Selbstadjungiertheit.
      \emph{(Skript [6].)}
  \end{enumerate}
\end{tcolorbox}

\begin{tcolorbox}[proofbox, title={Satz 6 — $\mathbb{Z}_3$-Restriktion und $\chi$-Twists \B}]
  \begin{enumerate}
    \item $[\Fhat, P_0] = 0$ ($\mathbb{Z}_3$-Kovarianz aus (F2))
      $\Rightarrow$ Restriktion auf den Orbifold-Sektor ist selbstadjungiert.
      \emph{(Skript [7].)}
    \item Fixpunkt-Randbedingungen $\chi^3 = 1$ (Dok.~322/320):
      unitär implementierte Twists; Spektrum $(n+\alpha)^2 \ge 0$ reell
      in allen drei Twist-Klassen.
      \emph{(Skript [8].)}
  \end{enumerate}
\end{tcolorbox}

\begin{tcolorbox}[proofbox, title={Korollar 7 — $\xi$ als wohldefiniertes Minimum \B}]
  $\Fhat_{D_4}$ (Restriktion auf $D_4$-Sektor, Dok.~314) ist selbstadjungiert.
  Nach dem Min-Max-Prinzip:
  \[
    \xi = \lambda_{\min}(\Fhat_{D_4})
        = \min_{\|\psi\|=1,\,\psi\in\Hcal_{D_4}} \langle\psi, \Fhat\psi\rangle
  \]
  \textbf{wohldefiniert und trunkierungsstabil}: die unterste Mode liegt
  in jedem Trunkierungsraum, der sie enthält.
  \emph{(Skript [9].)}
\end{tcolorbox}

% ============================================================
\section{Statusbilanz}

\begin{center}
\begin{tabular}{@{}lc@{}}
  \toprule
  \textbf{Aussage} & \textbf{Status} \\
  \midrule
  $L_nL_m = L_{\min}$ durch Filtrierung exakt realisiert & \B \\
  $\Fhat$ beschränkt durch endliche Stufenzahl ($\Lnull$) & \B \\
  Symmetrie als $\mathbb{Z}_3$-Folge, mit Gegenprobe & \B \\
  Selbstadjungiertheit, Defektindizes $(0,0)$, Eindeutigkeit & \B \\
  Fraktales Maß (endliche Rekursion) erhält Selbstadjungiertheit & \B \\
  $\mathbb{Z}_3$-Restriktion und alle drei $\chi$-Twist-Klassen s.a. & \B \\
  $\xi = \lambda_{\min}(\Fhat_{D_4})$ wohldefiniert, trunkierungsstabil & \B \\
  \midrule
  Restfälle & \textbf{keine} \\
  \bottomrule
\end{tabular}
\end{center}

Die R82-Lücke ist vollständig geschlossen.
Es gibt keine Unendlichkeiten in FFGFT;
$\sum r_n < \infty$ war nie eine Zusatzforderung,
sondern folgt direkt aus $\Lnull = \xi\cdot\lP$.

% ============================================================
\section{Vorgeschlagener Registereintrag R86}

\textbf{R86 -- Selbstadjungiertheit von $\hat{F}$ (August 2026):}
Dok.~327 schließt die in R82 als [S] deklarierte Lücke vollständig.
Die Endlichkeit der Skalenleiter ist durch $\Lnull = \xi\cdot\lP$
(Dok.~180 \K) garantiert — keine Unendlichkeiten in FFGFT,
damit $\Fhat$ beschränkt und $\Dcal(\Fhat) = \Hcal$.
Die Symmetrie $\Fhat = \Fhat^\dagger$ folgt aus der $\mathbb{Z}_3$-Paarung
$(k,-k)$ (dieselbe wie in Dok.~325 \B).
Defektindizes $(0,0)$: genau eine selbstadjungierte Realisierung \B.
Fraktales Maß (endliche 100er-Rekursion), $\mathbb{Z}_3$-Restriktion
und alle drei $\chi$-Twist-Klassen erhalten die Selbstadjungiertheit \B.
$\xi = \lambda_{\min}(\Fhat_{D_4})$ wohldefiniert und trunkierungsstabil \B.
Es verbleiben keine Restfälle.
Prüfskript: \texttt{327\_selbstadjungiertheit\_F\_verify.py}
(21 Assertionen, alle bestanden).\\
\emph{Deklariert 23.~August 2026}

% ============================================================
\section*{Bezug zum Korpus}

\begin{center}
\begin{tabular}{@{}ll@{}}
  \toprule
  \textbf{Baustein} & \textbf{Dok.} \\
  \midrule
  Minimale Längenskala $\Lnull = \xi\cdot\lP$ & 180 \\
  Definition $\Fhat$, $\Dcal(\Fhat)$, R82-Lücke & 322 \\
  $D_4$-Gitter im Hilbertraum & 314 \\
  Fixpunkt-Randbedingungen, Neutrino-Lokalisierung & 320, 322 \\
  $\mathbb{Z}_3$-Trialität, Sektorstruktur & 321 \\
  Sektorpaarung $(k,-k)$ — dieselbe wie im Hawking-Mechanismus & 325 \\
  Fraktales Maß, 100-fache Rekursion, $D_f = 3-\xi$ & 322 \\
  \bottomrule
\end{tabular}
\end{center}

\section*{Literatur}

\begin{enumerate}
  \item J.~Pascher, Dok.~180, 314, 320, 321, 322, 325, FFGFT-Korpus, 2026.
  \item M.~Reed, B.~Simon, \textit{Methods of Modern Mathematical Physics~I/II},
    Academic Press.
  \item J.~von~Neumann, \textit{Allgemeine Eigenwerttheorie Hermitescher
    Funktionaloperatoren}, Math.\ Ann.\ 102, 1929.
\end{enumerate}

\vspace{6pt}\noindent\small
\textcopyright\ 2026 Johann Pascher ·
\href{https://creativecommons.org/licenses/by/4.0/}{CC BY 4.0}
